\documentclass[12pt]{article}
\usepackage[margin=1in]{geometry}
\usepackage{amsmath,amssymb,amsthm}
\usepackage{mathtools}
\usepackage{hyperref}

\newtheorem{theorem}{Theorem}
\newtheorem{proposition}[theorem]{Proposition}
\newtheorem{corollary}[theorem]{Corollary}
\newtheorem{remark}[theorem]{Remark}
\newtheorem{definition}[theorem]{Definition}

\newcommand{\R}{\mathbb{R}}
\newcommand{\N}{\mathbb{N}}
\newcommand{\supp}{\operatorname{supp}}
\newcommand{\argmin}{\operatorname{arg\,min}}
\newcommand{\ol}{\overline}

\title{Endpoint 2 Beyond Finite $M$:\\
Countable-Atomic Conditional Extension}
\author{Formalization note prepared for Piotr Dworczak}
\date{March 17, 2026}

\begin{document}
\maketitle

\begin{abstract}
This note records the endpoint reached on the countable-atomic attainment route for extending the existence direction of Theorem~2 beyond finite message spaces.
The unconditional attainment theorem is false on this route because of escape of mass.
The surviving positive result is conditional:
if one near-optimal sublevel set is rowwise uniformly tight and the reduced value functional is lower semicontinuous in product-$\ell^1$, then the adversary problem attains a minimizer.
At that minimizer, the already-banked selector/subgradient argument yields rowwise equal-payoff-on-support, conditional on the explicit Bayes-side subgradient-realization caveat.
\end{abstract}

\tableofcontents

%% =====================================================================
\section{Summary}
%% =====================================================================

This endpoint has four parts.

\begin{enumerate}
\item We changed architectures. Instead of trying to recover exact recurring message fibers, we lifted the finite alternative proof to
\[
X := \prod_{\mu \in I} \ell^1(M),
\qquad
X^* := \prod_{\mu \in I} \ell^\infty(M),
\qquad
K := \prod_{\mu \in I} \Delta(M) \subset X,
\]
with $M$ countable and $I$ finite.

\item We proved that the unconditional attainment claim on this route is false.
There is a genuine escape-of-mass counterexample: minimizing sequences can run down an infinite message tail without converging in product-$\ell^1$.

\item We then identified the exact positive replacement.
If there exists $\eta > 0$ such that
\[
A_\eta := \{ \beta \in K : V(\beta) \le \inf_K V + \eta \}
\]
is rowwise uniformly tight, and if $V$ is lower semicontinuous, then $V$ attains its infimum on $K$.

\item Once a minimizer $\beta^*$ exists, the finite alternative-proof selector/subgradient mechanism survives on this branch.
With the explicit Bayes-side subgradient-realization caveat, one gets
\[
g^* \in \partial V(\beta^*),
\qquad
-g^* \in N_K(\beta^*),
\]
hence rowwise equal-payoff-on-support.
\end{enumerate}

The honest endpoint is therefore conditional.
We did not find a genuinely earlier model-side primitive that forces this tightness package.

%% =====================================================================
\section{Hand-Wavy Overview}
%% =====================================================================

In the finite proof, the hard step is not the first-order calculation itself.
The hard step is that the adversary actually has an optimizer.
When $M$ is finite, that is free because the strategy space is a compact finite-dimensional product of simplices.

When $M$ is countable, this breaks.
Mass can leak farther and farther into the message tail while the objective keeps improving.
So the first thing we proved is that unconditional attainment is simply false on this branch.

The right repair is compactness by tightness.
If all near-minimizers keep almost all of their mass on a common finite core of messages, then one recovers the missing precompactness in product-$\ell^1$.
That gives an actual minimizer.
Once the minimizer exists, the finite alternative proof becomes portable again:
subgradients plus the simplex normal cone produce the rowwise equal-payoff-on-support condition.

So the theorem does survive beyond finite $M$ on this route, but only conditionally:
we must explicitly forbid escape of mass.

%% =====================================================================
\section{Formal Setup}
%% =====================================================================

Fix a countable message space $M$ and a finite row index set $I$.
Write
\[
X := \prod_{\mu \in I} \ell^1(M),
\qquad
\|\beta\|_X := \sum_{\mu \in I} \|\beta_\mu\|_1,
\qquad
K := \prod_{\mu \in I} \Delta(M) \subset X.
\]

Let
\[
V : K \to \R
\]
be the reduced value functional on the countable-atomic attainment route.
Set
\[
c := \inf_{\beta \in K} V(\beta).
\]

For $\eta > 0$, define the near-optimal sublevel set
\[
A_\eta := \{ \beta \in K : V(\beta) \le c + \eta \}.
\]

\begin{definition}[Rowwise uniform tightness]
We say that a set $A \subset K$ is rowwise uniformly tight if for every $\varepsilon > 0$ and every row $\mu \in I$ there exists a finite set $F_{\mu,\varepsilon} \subset M$ such that
\[
\sup_{\beta \in A} \beta_\mu(M \setminus F_{\mu,\varepsilon}) \le \varepsilon .
\]
\end{definition}

%% =====================================================================
\section{Negative Result}
%% =====================================================================

\begin{proposition}[Unconditional attainment fails]
\label{prop:escape}
Under the standing hypotheses alone, the statement
\[
\inf_{\beta \in K} V(\beta)
\quad \text{is attained on } K
\]
is false on this route.
\end{proposition}

\begin{proof}
Take $M = \N$, let $I$ be any nonempty finite set, and define coefficients
\[
a_m := \frac{1}{m}.
\]
Consider the linear reduced functional
\[
V(\beta) := \sum_{\mu \in I} \sum_{m \ge 1} a_m \beta_\mu(m)
=
\sum_{\mu \in I} \sum_{m \ge 1} \frac{\beta_\mu(m)}{m}.
\]
This functional is bounded, convex, and continuous on $X$.

For each $n \ge 1$, define $\beta^n \in K$ by
\[
\beta^n_\mu := \delta_n
\qquad \text{for every } \mu \in I.
\]
Then
\[
V(\beta^n) = \frac{|I|}{n} \to 0,
\]
so $\inf_K V = 0$.
But for every $\beta \in K$, each row $\beta_\mu$ is a probability distribution and every coefficient $1/m$ is strictly positive, hence
\[
V(\beta) > 0.
\]
So the infimum is not attained.

This is exactly escape of mass: the sequence $(\beta^n)$ sends all mass to messages farther and farther out in the tail, and has no convergent subsequence in product-$\ell^1$.
\end{proof}

\begin{remark}
Proposition~\ref{prop:escape} is why endpoint 2 cannot be unconditional on the present record.
Any positive theorem on this branch must explicitly rule out this tail escape mechanism.
\end{remark}

%% =====================================================================
\section{Conditional Theorem}
%% =====================================================================

\begin{theorem}[Conditional attainment from one tight near-optimal sublevel]
\label{thm:attainment}
Assume $V$ is lower semicontinuous for the $X$-norm.
Suppose there exists $\eta > 0$ such that the near-optimal set
\[
A_\eta = \{ \beta \in K : V(\beta) \le c + \eta \}
\]
is rowwise uniformly tight.
Then there exists $\beta^* \in K$ such that
\[
V(\beta^*) = c = \inf_{\beta \in K} V(\beta).
\]
\end{theorem}

\begin{proof}
Choose a minimizing sequence $(\beta^n) \subset K$ satisfying
\[
V(\beta^n) \le c + \min\{\eta, 1/n\}.
\]
Then $\beta^n \in A_\eta$ for all $n$.

Because $I$ is finite, rowwise tightness can be merged into one product-tail bound:
for every $\varepsilon > 0$ there exists a finite set $F_\varepsilon \subset M$ such that
\[
\sup_{\beta \in A_\eta} \sum_{\mu \in I} \beta_\mu(M \setminus F_\varepsilon) \le \varepsilon.
\]
Indeed, choose finite rowwise truncation sets $F_{\mu,\varepsilon/|I|}$ and take their union.

Fix $\varepsilon_k = 2^{-k}$ and choose finite sets
\[
F_k \uparrow M
\]
such that
\[
\sup_{\beta \in A_\eta} \sum_{\mu \in I} \beta_\mu(M \setminus F_k) \le 2^{-k}.
\]

For each fixed $k$, the coordinate block
\[
\big(\beta_\mu(m)\big)_{(\mu,m)\in I\times F_k}
\]
ranges in a compact finite-dimensional box.
By diagonal extraction, after passing to a subsequence we may assume that for every $(\mu,m) \in I \times M$,
\[
\beta^n_\mu(m) \to \beta^*_\mu(m)
\]
for some nonnegative array $\beta^* = (\beta^*_\mu)_{\mu \in I}$.

We now upgrade coordinatewise convergence to convergence in $X$.
Fix $k$.
On the finite set $I \times F_k$,
\[
\sum_{\mu \in I}\sum_{m \in F_k} |\beta^n_\mu(m)-\beta^*_\mu(m)| \to 0.
\]
Also, by the choice of $F_k$,
\[
\sum_{\mu \in I}\beta^n_\mu(M \setminus F_k) \le 2^{-k}
\qquad\text{for all } n.
\]
Passing to the limit on $F_k$ shows
\[
\sum_{\mu \in I}\beta^*_\mu(M \setminus F_k) \le 2^{-k}.
\]
Hence
\[
\|\beta^n-\beta^*\|_X
\le
\sum_{\mu \in I}\sum_{m \in F_k} |\beta^n_\mu(m)-\beta^*_\mu(m)|
\;+\;
\sum_{\mu \in I}\beta^n_\mu(M \setminus F_k)
\;+\;
\sum_{\mu \in I}\beta^*_\mu(M \setminus F_k).
\]
First let $n \to \infty$, then let $k \to \infty$.
This yields
\[
\beta^n \to \beta^* \quad \text{in } X.
\]

In particular $\beta^* \in K$: each row is nonnegative, and the uniform tail bounds prevent loss of total mass.

Now use lower semicontinuity:
\[
V(\beta^*) \le \liminf_{n\to\infty} V(\beta^n) \le c.
\]
Since $c = \inf_K V$, we must have $V(\beta^*) = c$.
So the infimum is attained.
\end{proof}

%% =====================================================================
\section{Selector/Subgradient Corollary}
%% =====================================================================

\begin{corollary}[Endpoint 2]
\label{cor:endpoint2}
Assume the hypotheses of Theorem~\ref{thm:attainment}, and let $\beta^*$ be an attained minimizer.
Assume in addition the following Bayes-side realization property at $\beta^*$:

\medskip
\emph{every relevant subgradient $d \in \partial V(\beta^*)$ can be realized as the payoff array of a messagewise Bayes-optimal selector family.}

\medskip
Then there exists a messagewise Bayes-optimal selector family $g^* \in X^*$ such that
\[
g^* \in \partial V(\beta^*),
\qquad
-g^* \in N_K(\beta^*).
\]
Consequently, for every row $\mu \in I$, all messages in $\supp \beta^*(\cdot \mid \mu)$ yield the same selected payoff, and every off-support message yields a weakly higher selected payoff.
\end{corollary}

\begin{proof}
This is exactly the already-banked selector/subgradient proposition on the countable-atomic attainment route.
Theorem~\ref{thm:attainment} provides the missing minimizer $\beta^*$.
Once $\beta^*$ exists, the finite alternative-proof argument carries over in the Banach pair
\[
X = \prod_{\mu \in I} \ell^1(M),
\qquad
X^* = \prod_{\mu \in I} \ell^\infty(M).
\]
The normal cone to the countable simplex gives rowwise equal-payoff-on-support exactly as in the finite case.
\end{proof}

%% =====================================================================
\section{Frontier Status}
%% =====================================================================

The honest endpoint of this branch is now:

\begin{enumerate}
\item unconditional negative theorem:
generic attainment fails by escape of mass;
\item conditional positive theorem:
one tight near-optimal sublevel set plus lower semicontinuity implies attainment;
\item downstream conditional corollary:
attainment plus Bayes-side subgradient realization implies rowwise equal-payoff-on-support.
\end{enumerate}

What remains open is not another hidden gap inside this theorem package.
What remains open is whether one can derive the tightness hypothesis from a genuinely earlier
model-side primitive.
On the present proof record, no such earlier primitive has been found.

\end{document}
